\section{Midsummer Night Dream 2013}

Anyone attracted to political action should read those political scientists who describe how men actually are and act, not as they should in some imaginary world. Only then can men be led effectively. Most today will think that is immoral. To the contrary, the result is Order, Authority, and Unity … the highest morality you can and should expect in a state.

Man is faithful by nature; woman is promiscuous. Man wants to protect his own; this entails exclusion. Woman wants to nurture indiscriminately; this entails inclusion.

Those capable of thinking only in material terms and images or who understand time only in terms of number, will never grasp much of anything.

Eckhart said that the beginning of spirituality is dependent on recognizing what one is as a being outside time.

Identity is not biological. USA was 90\% white in my childhood, Boston much more, but that means nothing. In \textit{Virgin Spring}\footnote{\url{https://gornahoor.net/?p=1416}}, Töre was a noble, a father and patriarch, his daughter beautiful and blonde and virginal — they lived in the light. The maidservant was dark, promiscuous, and played with frogs. In our time, many think that is ``tradition" and are attracted to the bizarre, the abnormal, the ugly. The shepherds who gave into their craven lusts were also ethnic Swedes. So who then is a Swede according to the Identity movements? In Evola's expression, who is ``of race"? That is a spiritual quality and by the law of involution/evolution, not all who appear so are the bearers of the race.

Organizations eventually fail because the leaders become weak and indiscriminate; they have no standards or no longer adhere to them; the refuse to enforce discipline. Instead, they prefer to be ``nice" and ``inclusive", proving the hold that the modern mind still has on them. Stronger men know who is in and who is out. There was an order in the last century that banned non-fathers and atheists. In our time the new right and the identity right embrace them. They have their reasons, none of them good, and are more in the nature of confabulations.

Even a blog should strive to be what it describes, to the limited extent that is possible. That is why certain rules on commenting are established, more as an experiment. To reject a certain type of comment is not to reject all; that is the way boys think. One fellow who was asked to take some time out from commenting did so, but came back after a suitable time elapsed. That I can respect.

Androphiles try to act like men, so they use f**k a lot in their writings to show how macho they are. That may have had a chance to be true at one time, but now that even high school girls speak like that, it no longer has any impact. Do you want a handbook on acting manly? Learn to be chivalrous. I just saved you the price of a useless book.

Men don't relish considering themselves as victims.

Men are in tenuous positions today, forced into submission to those who have no right to authority over them. But even the slave Epictetus was free.

Thanks to facebook, I found out that there are certain things men are supposed to believe only because it is the 21st century. It saves all the trouble of thinking, using logic, and research; instead, it is only necessary to check one's calendar. Who knows what will be the cause du jour next year, in 5 years, in 10? Is there someone in charge of this or do all the locusts get hungry at the same time?

I don't care very much about people's opinions; what interests me is their inner state. What is it like to be them? What caste would they be, which spiritual race? It takes a few probing questions and to discern that by listening carefully to their answers. A man's knowledge is something different.

People create sand castles out of words and then try to live in them. Words have depth or perhaps not. Do they express something real or are they just nominal? Few can distinguish between the natural and the artificial.

I still remember my dreams as a baby and the first time I realized I was an ``I". But I haven't yet been able to remember who was dreaming the baby into existence.

Gurdjieff told the story about a magician who hypnotized his sheep into believing they were men. It did not work out well for the sheep because the charade was eventually found out and they felt worse than if they had simply accepted being sheep.

Esoteric writing is always incomplete, leaving something out, even being misleading, while expecting the correct readers to discern the secret meaning.

Who built the bus?



\flrightit{Posted on 2013-07-18 by Cologero }

\begin{center}* * *\end{center}

\begin{footnotesize}\begin{sffamily}



\texttt{Jason-Adam on 2013-07-19 at 10:17 said: }

Gurdjieff's story reminds me of some problems I've had with body-spirit conflicts in my life. It was hard to figure who is the ``real" me sometimes……..even know I still have problems centering around the Chief Feature.


\hfill

\texttt{Scardanelli on 2013-07-19 at 13:28 said: }

There was an evil magician. He lived deep in the mountains and the forests, and he had thousands of sheep. But the problem was that the sheep were afraid of the magician because every day the sheep were seeing that one of them was being killed for his breakfast, another was being killed for his lunch. So they ran away from the magician's ranch and it was a difficult job to find them in the vast forest. Being a magician, he used magic.

He hypnotized all the sheep and suggested to them first of all that they were immortal and that no harm was being done to them when they were skinned, that, on the contrary, it would be very good for them and even pleasant; secondly he suggested that the magician was a good master who loved his flock so much that he was ready to do anything in the world for them; and in the third place he suggested to them that if anything at all were going to happen to them it was not going to happen just then, at any rate not that day, and therefore they had no need to think about it.

He then told different sheep…to some, ``You are a man, you need not be afraid. It is only the sheep who are going to be killed and eaten, not you. You are a man just like I am." Some other sheep were told, ``You are a lion — only sheep are afraid. They escape, they are cowards. You are a lion; you would prefer to die than to run away. You don't belong to these sheep. So when they are killed it is not your problem. They are meant to be killed, but you are the most loved of my friends in this forest." In this way, he told every sheep different stories, and from the second day, the sheep stopped running away from the house.

They still saw other sheep being killed, butchered, but it was not their concern. Somebody was a lion, somebody was a tiger, somebody was a man, somebody was a magician and so forth. Nobody was a sheep except the one who was being killed. This way, without keeping servants, he managed thousands of sheep. They would go into the forest for their food, for their water, and they would come back home, believing always one thing: ``It is some sheep who is going to be killed, not you. You don't belong to the sheep. You are a lion — respected, honored, a friend of the great magician." The magician's problems were solved and the sheep never ran away again.'


\hfill

\texttt{scardanelli on 2013-07-19 at 13:46 said: }

From:

\url{http://glossary.cassiopaea.com/glossary.php?id=2}

The site has some useful clarifications of terms and concepts from Mouravieff and Gurdjieff, mixed in with a lot of other questionable material.


\hfill

\texttt{Avery Morrow on 2013-07-19 at 21:35 said: }

When the locusts get hungry they are trying to eat each other.


\hfill

\texttt{Cologero on 2013-07-22 at 13:00 said: }

That is part of the common problem, J-A. The ``I" is always there, it is not a ``me", and there is no figuring.


\hfill

\texttt{Cologero on 2013-07-22 at 13:02 said: }

AM, I assume that is the koan that your Zen Master gave you as you quit Japan … so when the locusts swarm, they are trying to fly away from each other to avoid being eaten.


\hfill

\texttt{CaseyAnn on 2013-10-03 at 02:01 said: }

The first part of this post has always stuck with me.

``Man is faithful by nature; woman is promiscuous. Man wants to protect his own; this entails exclusion. Woman wants to nurture indiscriminately; this entails inclusion."

I understand it in the sense of the second and third sentences; but I've always been troubled about it in the sense of the biological expression. Tonight Aristotle brought it up in my readings:

``[F]or the female is impregnated by one copulation, while the male impregnates many females."

Does it matter?


\hfill

\texttt{Synodius on 2013-10-03 at 13:20 said: }

I think Cologero meant man's attitude to logic (and rules) which is more binary, while a woman with her more integrative bias tends to make decisions based on emotions. For example polygamy as a biological fact does not go against `faithfulness'. only breaking of the rules that regulate it in a society does.


\end{sffamily}\end{footnotesize}
